\section{Closed-source image watermark
removal}\label{record-docs-closed-image-watermarking-md-closed-source-image-watermark-removal}

This experiment assumes that the attacker receives only generated RGB
images. It does not receive target-generator weights, decoder weights,
latent states, gradients, prompts, class labels, watermark keys,
payloads, detector weights or detector responses. It simulates this
access using the preserved PNG outputs of DistSeal RAR-XL and
UViT-H/DC-AE. It does not query a commercial API.

The selected method is \textbf{NFPA}, from \emph{The Future Unmarked:
Watermark Removal in AI-Generated Images via Next-Frame Prediction},
\href{https://proceedings.neurips.cc/paper_files/paper/2025/hash/b48405307576689bea14d221ccba73ae-Abstract-Conference.html}{NeurIPS
2025 main conference}. The full 26-page conference PDF, including
appendices, was downloaded and read. The
\href{https://github.com/1249748036/NFPA}{official implementation} is
pinned to \texttt{\seqsplit{7f449f305506f95876fbfde0690395bea720269c}}
and its pipeline is imported directly. The newer CVPR 2026 RAVEN
repository currently states that code/configurations are unreleased, so
it was not used as an approximate reproduction.

NFPA uses an independent public Stable Diffusion model to DDIM-invert an
input image. It searches bounded camera-motion flows that maximize the
difference of the public model's latent representation, then generates a
second frame using the authors' frame-attention mechanism. Neither the
search objective nor its stopping criterion queries a watermark
detector. This is output purification; it does not change the closed
provider's generator or its future outputs.

The released settings are retained: 10 inversion/denoising steps, empty
prompts, guidance 7.5, fp16 and xy motion. Delta 40 is the predeclared
primary setting; delta 20 provides a weaker setting and delta 0 provides
a regeneration control. Delta 0 still executes the released VAE/DDIM and
zero-flow grid resampling; it is not an unchanged-image control. A
separate resize-only arm measures the effect of Lanczos resizing to 512
and restoring the native output resolution. This restoration is needed
for the 256-pixel RAR outputs. DC-AE outputs are already 512 pixels. The
native-resolution image is always supplied to the owner.

The original \texttt{\seqsplit{stabilityai/stable-diffusion-2-1-base}}
endpoint returned HTTP 404 during setup. The independent attack model
therefore uses the public mirror
\texttt{\seqsplit{Manojb/stable-diffusion-2-1-base}}, revision
\texttt{\seqsplit{0094d483a120f3f33dafbd187ea4aa60d10de75c}}. Its local
fp16 weights are hashed in the task. This model-hosting difference is
disclosed; the experiment does not claim an independently verified
byte-for-byte match to the unavailable original host.

\subsection{Data and owner
evaluation}\label{record-docs-closed-image-watermarking-md-data-and-owner-evaluation}

For each target, the full experiment uses 512 marked images, 512
independent clean-generator controls and 256 independent clean
calibration images from the preserved image run. A preregistered early
pilot contains all 256 calibration images and 32 evenly spaced images
from each test cohort. Both zero-class and other-class images are
represented. The full run reuses the completed pilot images rather than
regenerating them.

Workers receive opaque image IDs and a public-model task. Separate
processes receive the private owner configuration. This is an enforced
input/argument contract, not an operating-system sandbox against a
malicious process sharing the filesystem account. Owner scoring occurs
after the complete pilot or full strength sweep has been sealed. Pilot
artifact indexes remain immutable when the full indexes grow. All
strengths are reported; owner scores never choose a strength, optical
flow, seed or retained output.

The owner strictly reloads the original detectors and first verifies
that saved unattacked images reproduce their existing bit scores.
Reports include bit accuracy, AUROC, detection and empirical false
positives with confidence intervals at both the original frozen
threshold and a threshold calibrated on the independently attacked
calibration images. Paired quality includes LPIPS, CLIP image-embedding
cosine, SSIM and PSNR. Pixel metrics are alignment-sensitive and can
penalize legitimate viewpoint changes, so they are diagnostics alongside
the perceptual metrics. None certifies that every object or fact was
preserved. The pilot's 32 clean test images and even the full 256
calibration images cannot precisely establish extremely low
false-positive rates. No large-sample FID reproduction is claimed.

\subsection{Execution}\label{record-docs-closed-image-watermarking-md-execution}

Prepared resources are stored under
\texttt{\seqsplit{results/research\_workflow/watermarking/resources/}}.
Experiments retain immutable source hashes and separate per-case image
tasks.

\begin{Verbatim}
PYTHONPATH=. uv run --no-sync python scripts/closed_image_watermark_experiments.py prepare
PYTHONPATH=. uv run --no-sync python scripts/closed_image_watermark_experiments.py submit RUN/manifest.json preflight
# After checking real-model timing and sample outputs:
PYTHONPATH=. uv run --no-sync python scripts/closed_image_watermark_experiments.py submit RUN/manifest.json
\end{Verbatim}

GPU preflight attacks one clean and one marked image per model, checks
native output dimensions, measures time per image and verifies owner
reloads. The two case lanes can run in parallel after preflight. Each
lane runs the pilot attack, pilot owner scoring, remaining full attack
and full owner scoring in separate processes. Jobs use one H100 and
two-hour dev QoS with advance signals and automatic requeue. Each
completed image and score batch is saved atomically. An interruption can
repeat the unfinished image/batch and startup work; zero compute loss is
not guaranteed. Runtime caches use project storage and are removed on
normal job exit. No environment is unpacked onto node temporary disk.

Focused validation:

\begin{Verbatim}
PYTHONPATH=. uv run --no-sync python -m pytest ai_content_detector/tests/test_closed_image_watermarking.py -q
\end{Verbatim}

The contract, integrity, geometry, pilot/full resumption and
worker-argument checks passed on CPU. The released NFPA module imports
in the existing main environment. GPU measurements belong to each run's
preflight and result files; preparation and submission are not completed
attack results.

\subsection{Validated startup and submitted
evaluation}\label{record-docs-closed-image-watermarking-md-validated-startup-and-submitted-evaluation}

Frozen run: \texttt{\seqsplit{20261104T082631753968Z}}. GPU preflight
\texttt{\seqsplit{1939394}} completed on an H100 80GB in 3m23s. Both
owner detectors exactly reproduced preserved clean and marked image
scores. The released attack generated the expected native output sizes.
Warm attack times were 0.70s/image for RAR and 0.74s/image for DC-AE;
RAR's first cold image took 5.63s. Peak allocated GPU memory was 3.81
GiB. This four-image timing check is a startup measurement, not a broad
benchmark.

The predeclared clean/marked visual pairs are saved in
\texttt{\seqsplit{preflight\_pairs.png}}. They preserve the broad scene
but show changes to faces, hands, fish details and framing. These
visible changes are a quality limitation to quantify alongside removal;
they must not be described as exact semantic or identity preservation.
The weaker strength and resize/regeneration controls remain in the full
report. No watermark removal success rate has yet been measured by this
preflight.

Evaluation lanes: RAR \texttt{\seqsplit{1939549}}, DC-AE
\texttt{\seqsplit{1939550}}, both dependent on successful preflight
\texttt{\seqsplit{1939394}}. Each lane retains its 32-image test pilot
and then extends to 512 test images per cohort without regenerating
completed outputs. Results are written to
\texttt{\seqsplit{model01/results\_pilot.json}} and
\texttt{\seqsplit{model02/results\_pilot.json}}, then the corresponding
\texttt{\seqsplit{results.json}} files. Pending jobs are not results.

Both full evaluation jobs completed successfully. See the
\href{records/docs/closed_image_results.tex}{completed results} for
detection, false positives and paired quality at all fixed strengths.
